\documentclass[a4paper,12pt,openany]{report}
\usepackage{eso-pic}
\usepackage{tikz}
\usetikzlibrary{shadows.blur}
\usepackage{booktabs}  % برای جدول زیباتر
\usepackage{colortbl}  % رنگ زدن ردیف‌ها
\usepackage{multirow}  % ادغام سلول‌ها برای سلسله مراتب
\usepackage{enumitem}
\usepackage{pifont}
\usepackage{graphicx}
\usepackage{calc} % برای محاسبه طول متن
\usepackage{fancyhdr}
\usepackage{array}
\usepackage[font=bf]{caption}
\usepackage{titlesec}
\usepackage{tcolorbox}
\usepackage{xepersian}
\usepackage{xcolor}
% -----------------------------------------------------------------
% این بخش شامل تنظیمات سبک نگارش و ظاهر مستند است
% نباید تغییر کند
%------------------------------------------------------------------
% تنظیم فونت فارسی
\settextfont[
  Path=fonts/,
  UprightFont=BNazanin.ttf,
  BoldFont=BNazaninBold.ttf
]{BNazanin}
\renewcommand{\baselinestretch}{1.5} % فاصله خطوط 1.5 برابر حالت معمول
\newfontfamily\symbolfont{DejaVu Sans}
\setlist[itemize]{label={\symbolfont\textbullet}, left=1.0em}

\newfontfamily\persiantitlefont[
  Path=fonts/,
  UprightFont=BTitr.ttf
]{BTitr}

% متن های لاتین را کوچکتر میکند تا در کنار حروف فارسی نمایش درستی داشته باشند
\setlatintextfont[Scale=0.85]{Roboto}

\makeatletter
\renewcommand{\footnoterule}{%
  \noindent
  \color{blue!60!black}
  \begin{flushright}
    \rule{0.3\columnwidth}{0.4pt}
  \end{flushright}
  \kern 4pt
}
\makeatother

\usepackage[bottom]{footmisc}
\usepackage{setspace}
\renewcommand{\footnotelayout}{\normalfont\color{black}\raggedleft\setstretch{0.9}}

% تنظیم استایل فصل ها
\titleformat{\chapter}[display]
  {\normalfont\huge\bfseries}
  {\chaptername\ \thechapter}{10pt}{\Huge}

\titlespacing*{\chapter}{0pt}{-20pt}{20pt}

% تنظیم استایل بخش‌ها
\titleformat{\section}
  {\normalfont\Large\bfseries\color{headercolor}} % فونت بزرگ، پررنگ، رنگ آبی
  {\thesection}{1em}{} % شماره و فاصله

\titlespacing*{\section}
  {0pt}{3ex plus 1ex minus .2ex}{2ex plus .2ex} % فاصله قبل و بعد

% تنظیم استایل زیربخش
\titleformat{\subsection}
  {\normalfont\large\bfseries\color{orgblue}} % کمی کوچکتر، پررنگ، آبی تیره
  {\thesubsection}{1em}{}

\titlespacing*{\subsection}
  {0pt}{2.5ex plus 1ex minus .2ex}{1.5ex plus .2ex}

% تنظیم استایل زیرزیربخش با رنگ خاکستری تیره برای خوانایی روی بک‌گراند سفید
\titleformat{\subsubsection}
  {\normalfont\normalsize\bfseries\color{gray!30!black}} % خاکستری تیره
  {\thesubsubsection}{1em}{}

\titlespacing*{\subsubsection}
  {0pt}{2ex plus 1ex minus .2ex}{1ex plus .2ex}



\begin{document}

\thispagestyle{empty} % بدون شماره صفحه

\begin{center}

\vspace*{0.2cm}

\includegraphics[width=6cm]{iust_logo.png} % لوگوی دانشگاه (تغییر نام فایل اگر لازم است)

\vspace{2cm}

{\LARGE \textbf{دانشگاه علم و صنعت ایران}\par}

\vspace{1cm}

{\large دانشکده [نام دانشکده، مثلاً مهندسی برق]\par}

\vspace{0.5cm}

{\large گروه آموزشی [نام گروه، مثلاً مخابرات]\par}

\vspace{3cm}

{\huge \textbf{گزارش پیشرفت رساله دکتری}\par}

\vspace{2cm}

{\Large در رشته: مهندسی [نام رشته] - [گرایش]\par}

\vspace{2cm}

{\LARGE تحلیل تطبیق‌پذیر ترافیک شبکه در محیط‌های رانشی \par}

\vspace{3cm}

{\Large محمدجواد فقیه‌نیا \par}
{\large شماره دانشجویی\par}

\vspace{2cm}

{\Large استادان راهنما\par}
\vspace{0.5cm}
{\Large دکتر رضا انتظاری ملکی و دکتر سعید پارسا \par}

\vspace{1cm}


\vspace{0.5cm}
{\Large دی ۱۴۰۴ \par}

\vspace{2cm}


\end{center}

\clearpage

% اینجا می‌توانید بقیه گزارش را اضافه کنید

\end{document}